\section{The model}\label{sec:model}

This section fixes the market, the two optimisation problems, the core, the scaling and the
comparison model. A table of notation is at the end.

\subsection{The market}\label{sec:model-market}

There are three sectors $A$, $B$, $C$, each with $q$ agents, written $a_i$, $b_j$, $c_k$ for
$1\le i,j,k\le q$; when no index is needed we write $a$, $b$, $c$ for agents of the three
sectors. A \emph{pair} is a set of two agents from different sectors. There are $3q^2$ pairs.
Each pair $p$ has a \emph{value} $w(p)$, and the $3q^2$ values are independent and uniform on
$[0,1]$. The \emph{deficit} of a pair is
\begin{equation}\label{eq:model-pairdeficit}
\defi(p)=1-w(p),
\end{equation}
so the pair deficits are again independent and uniform on $[0,1]$.

A \emph{team} is a triple $e=(a,b,c)$ with one agent from each sector. It contains the three
pairs $\{a,b\}$, $\{b,c\}$, $\{a,c\}$. Its \emph{value} and its \emph{deficit} are
\begin{equation}\label{eq:model-teamdeficit}
\begin{aligned}
v(e)&=w(a,b)+w(b,c)+w(a,c),\\
\Defi(e)&=\defi(a,b)+\defi(b,c)+\defi(a,c)=3-v(e).
\end{aligned}
\end{equation}
There are $q^3$ teams. Two teams that share a pair have dependent values; two teams that share
no pair have independent values. Since $\Defi(e)$ is a sum of three independent uniform
variables,
\begin{equation}\label{eq:model-oneteam}
\P(\Defi(e)\le t)=\frac{t^3}{6}\qquad(0\le t\le1),
\end{equation}
and the density of $\Defi(e)$ is $t^2/2$ on $[0,1]$ and at most $t^2/2$ on $[0,3]$.

\subsection{Assignments and the fractional relaxation}\label{sec:model-lp}

A \emph{perfect assignment} is a set of $q$ pairwise disjoint teams; it covers every agent
exactly once. Its value is the sum of the values of its teams. We write
\begin{equation}\label{eq:model-T}
\Tq=\max_N\sum_{e\in N}v(e),\qquad
\Phiq=3q-\Tq=\min_N\sum_{e\in N}\Defi(e),
\end{equation}
where $N$ ranges over the perfect assignments.
A perfect assignment is determined by its set of $3q$ pairs, so two different assignments have
equal value only on a null set, and the maximiser is almost surely unique. We call it the
\emph{optimal assignment} and write $\Mopt$ for it.

A \emph{fractional assignment} is a vector $x=(x_e)$ of nonnegative weights on the teams with
\begin{equation}\label{eq:model-load}
\sum_{e\ni i}x_e=1\qquad\text{for every agent } i .
\end{equation}
The \emph{fractional relaxation} is the linear programme
\begin{equation}\label{eq:model-LP}
\LPq=\max_x\sum_e v(e)\,x_e,\qquad
\PhiLP=3q-\LPq=\min_x\sum_e\Defi(e)\,x_e ,
\end{equation}
where $x$ ranges over the fractional assignments and the second form uses $\sum_e x_e=q$. The indicator vector of a perfect assignment is a
fractional assignment, so $\LPq\ge\Tq$. The \emph{gap} is
\begin{equation}\label{eq:model-gap}
\Dq=\LPq-\Tq=\Phiq-\PhiLP\ \ge\ 0 .
\end{equation}
The \emph{support} of a fractional assignment $x$ is the set of teams with $x_e>0$.

The dual of \eqref{eq:model-LP} has one variable $u_i$ for each agent, unrestricted in sign
because the constraints \eqref{eq:model-load} are equations. Writing $u(e)=u_a+u_b+u_c$ for
$e=(a,b,c)$, it is
\begin{equation}\label{eq:model-dual}
\min\Bigl\{\sum_i u_i:\ u(e)\ge v(e)\text{ for every team } e\Bigr\}.
\end{equation}
The programme \eqref{eq:model-LP} is feasible and bounded. By the duality theorem of linear
programming the minimum in \eqref{eq:model-dual} is attained and equals $\LPq$, and for every
optimal $x$ and every optimal $u$,
\begin{equation}\label{eq:model-cs}
x_e>0\quad\Longrightarrow\quad u(e)=v(e).
\end{equation}

\subsection{The core}\label{sec:model-core}

The market defines a cooperative game with transferable utility on the set $I$ of the $3q$
agents. A coalition $S\subseteq I$ can form disjoint teams from its own members, and its worth
$v(S)$ is the largest total value of a set of pairwise disjoint teams contained in $S$ (zero if
$S$ contains no team). In particular $v(I)=\Tq$, because all values are nonnegative and the
sectors have equal size, so that a set of disjoint teams extends to a perfect assignment
without loss of value.

\begin{definition}\label{def:model-core}
The \emph{core} is the set of payoff vectors $p\in\R^I$ with $\sum_{i\in I}p_i=v(I)$ and
$\sum_{i\in S}p_i\ge v(S)$ for every coalition $S\subseteq I$.
\end{definition}

\begin{lemma}\label{lem:model-coreform}
A vector $p\in\R^I$ is in the core if and only if
\begin{equation}\label{eq:model-coreform}
p_i\ge0\ \text{ for every agent } i,\quad p(e)\ge v(e)\ \text{ for every team } e,\quad
\sum_{i\in I}p_i=\Tq ,
\end{equation}
where $p(e)=p_a+p_b+p_c$ for $e=(a,b,c)$.
\end{lemma}

\begin{proof}
If $p$ is in the core, apply the definition to $S=\{i\}$, whose worth is $0$, to
$S=\{a,b,c\}$, whose worth is $v(e)$, and to $S=I$. Conversely, let $p$ satisfy
\eqref{eq:model-coreform} and let $S$ be a coalition. Choose disjoint teams
$e_1,\dots,e_k\subseteq S$ with $v(S)=\sum_j v(e_j)$. Since $p\ge0$,
$\sum_{i\in S}p_i\ge\sum_j p(e_j)\ge\sum_j v(e_j)=v(S)$.
\end{proof}

\begin{proposition}\label{prop:model-core}
The core is nonempty if and only if $\Dq=0$.
\end{proposition}

\begin{proof}
Let $p$ be in the core. By Lemma~\ref{lem:model-coreform}, $p$ is feasible for the dual
\eqref{eq:model-dual} and has total $\Tq$. By weak duality $\LPq\le\sum_i p_i=\Tq\le\LPq$, so
$\Dq=0$.

Conversely, let $\Dq=0$ and let $u$ be an optimal solution of \eqref{eq:model-dual}, so that
$\sum_i u_i=\LPq=\Tq$ and $u(e)\ge v(e)\ge0$ for every team. The vector $u$ may have negative
entries, and we shift it sector by sector. Put $m_A=\min_{a\in A}u_a$ and define $m_B$, $m_C$ in
the same way. Every triple is a team, so $s:=m_A+m_B+m_C=\min_e u(e)\ge0$. Define
\[
p_a=u_a-m_A+\tfrac{s}{3},\qquad p_b=u_b-m_B+\tfrac{s}{3},\qquad p_c=u_c-m_C+\tfrac{s}{3}.
\]
Then $p_i\ge s/3\ge0$ for every agent, $p(e)=u(e)-s+s=u(e)\ge v(e)$ for every team, and, since
each sector has $q$ agents, $\sum_i p_i=\sum_i u_i-qs+qs=\Tq$. By
Lemma~\ref{lem:model-coreform}, $p$ is in the core.
\end{proof}

Thus the statement ``the core is empty'' is the statement $\Dq>0$: the optimal assignment is
not an optimal solution of the fractional relaxation. From here on the game plays no role, and
the paper is about the random variable $\Dq$.

\subsection{Deficits and the scaling}\label{sec:model-scaling}

A fixed agent lies in $q^2$ teams. By \eqref{eq:model-oneteam} the expected number of these
with deficit at most $t$ is $q^2t^3/6$ for $t\le1$, which is of order one when $t$ is of order
$q^{-2/3}$. We therefore measure team deficits on that scale and define the \emph{rescaled
cost}
\begin{equation}\label{eq:model-cost}
\cost_e=q^{2/3}\,\Defi(e).
\end{equation}
For $K\le q^{2/3}$ the expected number of teams $e$ at a given agent with $\cost_e\le K$ is
exactly $K^3/6=\Lam([0,K])$, where
\begin{equation}\label{eq:model-Lam}
\Lam(\dd\xi)=\frac{\xi^2}{2}\dd\xi\qquad\text{on }[0,\infty).
\end{equation}
A perfect assignment whose teams have rescaled costs of order one has total deficit of order
$q\cdot q^{-2/3}=q^{1/3}$. This is the order of $\Phiq$: by \cite{FPT20}, Theorem~2,
$\E\Phiq$ lies between two constant multiples of $q^{1/3}$. All the constants of this paper
are constants in front of $q^{1/3}$:
\[
\frac{\Phiq}{q^{1/3}}=\frac1q\sum_{e\in\Mopt}\cost_e ,
\qquad
\frac{\PhiLP}{q^{1/3}}=\frac1q\sum_e\cost_e\,x_e\quad\text{for an optimal } x .
\]
Each is an average cost per team, equivalently per agent of one sector.

For $K>0$ the \emph{near-tight hypergraph} $H_q(K)$ has the $3q$ agents as vertices and the
teams with $\cost_e\le K$ as hyperedges, each marked by its cost $\cost_e$. Its local limit is
the Poisson hypertree $\tree$ of Section~\ref{sec:tree} restricted to teams of cost at most
$K$. In $\tree$ every agent has countably many teams whose costs form a Poisson process of
intensity $\Lam$, every team has three agents, and the incidence graph is a tree; after the
restriction an agent has a Poisson number of teams with mean $K^3/6$. The precise statement
is Theorem~\ref{thm:localconv}.

\subsection{The independent-team model}\label{sec:model-iid}

In the market, teams that share a pair have dependent deficits. The \emph{independent-team
model} removes this dependence and keeps everything else. It has the same $3q$ agents and
$q^3$ teams, and the team deficits $\Defi(e)$ are independent, each distributed as the sum of
three independent uniform variables on $[0,1]$. So a single team has the same law
\eqref{eq:model-oneteam} in both models, and there are no pair values. Perfect assignments,
rescaled costs and the near-tight hypergraph are defined as in the market. Where both models
appear together, $\P$ and $\E$ refer to the market and $\P'$ and $\E'$ to the independent-team
model.

The integer lower bound is first proved in the independent-team model (Theorem~\ref{thm:iid})
and then carried to the market by a comparison inequality (Theorem~\ref{thm:transfer}). In the
market, a team with $\cost_e\le K$ shares a pair with about $\tfrac{9}{10}K^2q^{-1/3}$ other
such teams on average. This tends to zero, but slowly, and it is the reason why the two models
agree only at very large $q$ (Section~\ref{sec:discussion}).

\subsection{Notation}\label{sec:model-notation}

Table~\ref{tab:model-notation} lists the symbols used throughout. Logarithms are natural.
``With high probability'' means with probability tending to $1$ as $q\to\infty$.

\begin{table}[ht]
\caption{Notation.}\label{tab:model-notation}
\small
\begin{tabular}{@{}ll@{}}
\toprule
Symbol & Meaning \\
\midrule
$A$, $B$, $C$; $a_i$, $b_j$, $c_k$ & the three sectors; their agents \\
$w(p)$, $\defi(p)=1-w(p)$ & value and deficit of a pair $p$ \\
$v(e)$, $\Defi(e)=3-v(e)$ & value and deficit of a team $e$ \\
$\Tq$ & best total value of a perfect assignment \\
$\Mopt$ & the optimal assignment (almost surely unique) \\
$\Phiq=3q-\Tq$ & integer deficit \\
$\LPq$ & optimum of the fractional relaxation \\
$\PhiLP=3q-\LPq$ & fractional deficit \\
$\Dq=\LPq-\Tq$ & the gap; the core is nonempty iff $\Dq=0$ \\
$u$ & a solution of the dual \eqref{eq:model-dual} \\
$\cost_e=q^{2/3}\Defi(e)$ & rescaled cost of a team \\
$H_q(K)$ & near-tight hypergraph: teams with $\cost_e\le K$ \\
$\P'$, $\E'$ & probability, expectation in the independent-team model \\
$\Lam$ & the measure $\tfrac12\xi^2\dd\xi$ on $[0,\infty)$ \\
$\tree$ & the Poisson hypertree (Section~\ref{sec:tree}) \\
$\RDE$ & $\RDE(\mu)$ is the law of $\min_k(\xi_k-X_k'-X_k'')$; see \eqref{eq:intro-rde} \\
$\disp$ & dispersive order (Section~\ref{sec:existence}) \\
IFR & increasing failure rate \\
$\betaint=2.302$ & integer constant (Theorems~\ref{thm:iid}, \ref{thm:integer}) \\
$\alphalp=2.2519$ & bound on the tree cost (Theorem~\ref{thm:existence}) \\
$c_2$ & the constant of \cite{FPT20}, Theorem~2(b) \\
$C_0$ & $\sup_n\E\Phi_n/n^{1/3}$ \\
\bottomrule
\end{tabular}
\end{table}
